\documentclass[11pt, a4paper]{article}
\usepackage[a4paper, margin=2.2cm]{geometry}
\usepackage[utf8]{inputenc}
\usepackage[T1]{fontenc}
\usepackage[spanish]{babel}
\usepackage{lmodern}
\usepackage{enumitem}\setlist[enumerate]{topsep=3pt, itemsep=2pt}
\usepackage{booktabs}
\usepackage{tabularx}
\usepackage{xcolor}
\usepackage{hyperref}

\definecolor{primaryred}{HTML}{B91C1C}
\definecolor{darkslate}{HTML}{1E293B}
\definecolor{mutedgray}{HTML}{475569}
\definecolor{lightbg}{HTML}{F8FAFC}

\hypersetup{
    colorlinks=true,
    linkcolor=primaryred,
    urlcolor=primaryred,
    pdftitle={Sistema Multiagente para Análisis Inteligente de Expedientes de Inmuebles}
}

\setlength{\parindent}{0pt}
\setlength{\parskip}{5pt}

\begin{document}

% Encabezado del documento
\begin{center}
    {\Large \textbf{\color{darkslate} Sistema Multiagente para Análisis Inteligente de Expedientes de Inmuebles}}\\[6pt]
    {\normalsize \textbf{\color{primaryred} Hackatón Tiendas 3B $\cdot$ Reto 1 $\cdot$ Expedientes}}\\[4pt]
    {\small \color{mutedgray} \textbf{Equipo:} [por definir] \quad $\cdot$ \quad \textbf{Fecha de entrega:} [por definir] \quad $\cdot$ \quad \textbf{Estado:} en producción}
\end{center}

\vspace{4pt}
\hrule
\vspace{8pt}

% Índice
\tableofcontents

\vspace{10pt}
\noindent\fbox{\parbox{\dimexpr\textwidth-2\fboxsep-2\fboxrule\relax}{
    \small \textbf{Nota:} los datos marcados con \texttt{[por definir]} o \texttt{[por confirmar]} dependen de información que el comité organizador entrega al arranque del reto (criterios y reglas definitivos por tipo documental) o de decisiones pendientes del equipo. Ante cualquier diferencia, mandan los documentos oficiales del comité.
}}
\vspace{6pt}

% ---------------------------------------------------------
\section{Estructura del equipo y asignación de roles}
% ---------------------------------------------------------

Nos organizamos bajo SCRUM, con sus tres roles oficiales: Scrum Master, Product Owner y Development Team. Cada integrante tiene un área de enfoque principal, pero en la práctica todos colaboramos según las necesidades de cada sprint.

\begin{itemize}
    \item \textbf{Scrum Master — [Nombre]:} Coordina al equipo, da seguimiento a tareas y facilita los eventos SCRUM. Gestiona el tablero de tareas y cuida que el equipo trabaje sin bloqueos, sobre todo por el poco tiempo de un hackatón.
    \item \textbf{Product Owner / UX Research — [Nombre]:} Representa las necesidades del equipo legal de Tiendas 3B. Levanta requerimientos, redacta las historias de usuario, prioriza el backlog y valida que el prototipo cumpla con lo esperado. Prepara, junto con QA, el pitch de 3 minutos.
\end{itemize}

\subsection*{Development Team}

\begin{itemize}
    \item \textbf{Líder técnico / Arquitecto de agentes — [Nombre]:} Define la arquitectura multiagente, el diagrama de agentes, los patrones de diseño y el esquema de estado compartido. Gestiona el repositorio del equipo.
    \item \textbf{Desarrollador de agentes (LangChain) — [Nombre]:} Implementa los agentes, sus prompts, herramientas y salidas estructuradas: clasificador, extractores, validadores, agente de validación cruzada, verificador y dictaminador.
    \item \textbf{Desarrollador backend e ingesta — [Nombre]:} Construye la ingesta de PDF (texto y OCR), la segmentación con referencias de página, la API del análisis y la persistencia de resultados.
    \item \textbf{Desarrollador frontend — [Nombre]:} Desarrolla la interfaz del prototipo: carga de expediente, ejecución del análisis y presentación de resultados, hallazgos y evidencia.
    \item \textbf{QA / Documentación — [Nombre]:} Diseña y ejecuta las pruebas sobre expedientes de prueba contra el dictamen de referencia, elabora los diagramas, lleva la bitácora y documenta aspectos de seguridad y confidencialidad.
\end{itemize}

% ---------------------------------------------------------
\section{Descripción del proyecto}
% ---------------------------------------------------------

\subsection*{Contexto}
Uno de los principales retos para el crecimiento de Tiendas 3B es identificar y habilitar locales comerciales e industriales que cumplan los requisitos para operar como nuevas tiendas o centros de distribución. De cada inmueble se recibe documentación legal y administrativa que el equipo legal revisa antes de decidir si el expediente continúa.

\subsection*{Problema identificado}
Buena parte de la revisión es manual: leer documentos, localizar información relevante, validar requisitos e identificar problemas. Esto genera tres problemas principales:
\begin{itemize}
    \item Tiempo de personal especializado consumido en revisiones iniciales repetitivas.
    \item Límite en el número de expedientes que se pueden procesar.
    \item Riesgo de pasar por alto inconsistencias entre documentos (direcciones, propietarios, superficies, números oficiales) o de descartar expedientes que sí podían continuar.
\end{itemize}

\subsection*{Solución propuesta}
Un prototipo funcional basado en una arquitectura multiagente sobre LangChain que recibe un expediente, clasifica sus documentos, extrae los datos relevantes de cada tipo documental, los valida contra los criterios y reglas del comité, realiza validaciones cruzadas entre documentos del mismo inmueble y emite una recomendación inicial (\textit{Cumple}, \textit{Requiere revisión} o \textit{No cumple}) con su nivel de confianza y la evidencia que sustenta cada hallazgo.

El objetivo no es sustituir la decisión del equipo legal, sino permitirle revisar más expedientes en menos tiempo, detectar pronto los que claramente no cumplen y concentrar la intervención humana donde sí se necesita.

\subsection*{Alcance}
\begin{itemize}
    \item \textbf{Incluye:} los cuatro tipos documentales mínimos (Escrituras, Predial, Alineamiento y Número Oficial, Uso de Suelo), el análisis de punta a punta de un expediente, la interfaz del prototipo y el pitch.
    \item \textbf{No incluye:} la decisión legal final, el Reto 2 (Validat3B) ni la integración con sistemas internos de Tiendas 3B.
\end{itemize}

\subsection*{Propuesta de valor}
\begin{itemize}
    \item Reduce el tiempo de la revisión inicial y permite procesar más expedientes.
    \item Cada hallazgo cita el documento, la página y el fragmento que lo sustenta, de modo que el equipo legal entiende la recomendación sin abrir el expediente completo.
    \item Ante la duda, el sistema recomienda revisión humana en lugar de descartar, para no perder expedientes viables.
    \item Los documentos se tratan con confidencialidad: se usan solo dentro del reto y no se comparten fuera del equipo.
\end{itemize}

% ---------------------------------------------------------
\section{Formalización de los requerimientos funcionales}
% ---------------------------------------------------------

Los requerimientos funcionales describen qué debe hacer el sistema. Se levantaron a partir del planteamiento del Reto 1 y de los cuatro paquetes de documentos de ejemplo (12 de Escrituras, 16 de Predial, 18 de Alineamiento y Número Oficial y 16 de Uso de Suelo). Los criterios y reglas definitivos se incorporarán cuando el comité los entregue.

\subsection*{Gestión del expediente}

\textbf{RF-01 Carga de expediente}
\begin{itemize}
    \item \textbf{Descripción:} el sistema deberá permitir al usuario ingresar un expediente compuesto por uno o varios documentos PDF asociados a un mismo inmueble.
    \item \textbf{Razones:} es el punto de entrada del análisis; permite procesar expedientes completos y no documentos aislados.
    \item \textbf{Dependencias:} ninguna.
    \item \textbf{Fuentes:} planteamiento del Reto 1 (``permita ingresar un expediente'').
\end{itemize}

\textbf{RF-02 Identificación y clasificación de documentos}
\begin{itemize}
    \item \textbf{Descripción:} el sistema deberá identificar y clasificar automáticamente el tipo de cada documento recibido (Escrituras, Predial, Alineamiento y Número Oficial, Uso de Suelo u otro/no reconocido).
    \item \textbf{Razones:} elimina la clasificación manual y permite enrutar cada documento al agente especializado correcto.
    \item \textbf{Dependencias:} RF-01.
    \item \textbf{Fuentes:} planteamiento del Reto 1.
\end{itemize}

\textbf{RF-03 Documentos recibidos y faltantes}
\begin{itemize}
    \item \textbf{Descripción:} el sistema deberá determinar qué tipos documentales mínimos están presentes en el expediente y cuáles faltan.
    \item \textbf{Razones:} un expediente incompleto no puede dictaminarse como Cumple; el faltante debe ser visible desde el inicio.
    \item \textbf{Dependencias:} RF-02.
    \item \textbf{Fuentes:} planteamiento del Reto 1; guía por tipo documental del comité.
\end{itemize}

\subsection*{Extracción y trazabilidad}

\textbf{RF-04 Extracción de información por tipo documental}
\begin{itemize}
    \item \textbf{Descripción:} el sistema deberá extraer los datos relevantes definidos para cada tipo documental.
    \item \textbf{Razones:} es la base de las validaciones por documento y cruzadas.
    \item \textbf{Dependencias:} RF-02.
    \item \textbf{Fuentes:} guía por tipo del comité. Datos mínimos, según el planteamiento:
    \begin{itemize}
        \item \textbf{Escrituras:} titular, dirección e identificación del inmueble, superficie del terreno, datos relevantes y datos que lo relacionan con el resto del expediente.
        \item \textbf{Predial:} propietario o contribuyente, dirección, cuenta o clave predial, superficie registrada.
        \item \textbf{Alineamiento y Número Oficial:} dirección oficial, número oficial, identificación del predio, alineamiento determinado, fecha de emisión y vigencia, autoridad emisora.
        \item \textbf{Uso de Suelo:} dirección e identificación del inmueble, uso autorizado, compatibilidad con el uso comercial requerido, restricciones o condicionantes, superficie aplicable, fecha de emisión y vigencia, autoridad emisora.
    \end{itemize}
\end{itemize}

\textbf{RF-05 Evidencia por cada dato y hallazgo}
\begin{itemize}
    \item \textbf{Descripción:} el sistema deberá asociar cada dato extraído y cada hallazgo con su evidencia: documento, página y fragmento original que lo sustenta.
    \item \textbf{Razones:} las conclusiones deben ser trazables (regla del reto) y auditables por el equipo legal.
    \item \textbf{Dependencias:} RF-04.
    \item \textbf{Fuentes:} criterio de Explicabilidad y trazabilidad (20\%).
\end{itemize}

\subsection*{Validación}

\textbf{RF-06 Validación contra criterios y reglas}
\begin{itemize}
    \item \textbf{Descripción:} el sistema deberá validar la información extraída de cada documento contra los criterios y reglas proporcionados para su tipo documental, indicando requisitos cumplidos y no cumplidos.
    \item \textbf{Razones:} automatiza la verificación de requisitos de la revisión inicial.
    \item \textbf{Dependencias:} RF-04, RF-05.
    \item \textbf{Fuentes:} reglas oficiales del comité.
\end{itemize}

\textbf{RF-07 Detección de incompletitud, riesgos y faltantes}
\begin{itemize}
    \item \textbf{Descripción:} el sistema deberá detectar documentos incompletos, ilegibles o vencidos, información faltante y riesgos (por ejemplo, restricciones o condicionantes de uso).
    \item \textbf{Razones:} permite identificar pronto problemas que bloquean la continuidad.
    \item \textbf{Dependencias:} RF-06.
    \item \textbf{Fuentes:} planteamiento del Reto 1.
\end{itemize}

\textbf{RF-08 Validaciones cruzadas entre documentos}
\begin{itemize}
    \item \textbf{Descripción:} el sistema deberá contrastar entre documentos del mismo inmueble la dirección, el propietario o titular, la superficie, el número oficial y la identificación del predio, y detectar documentos que posiblemente correspondan a inmuebles diferentes. Para cada discrepancia deberá mostrar qué documentos la presentan, qué valor tiene cada uno y dónde está la evidencia.
    \item \textbf{Razones:} es la principal fuente de errores no detectados en la revisión manual.
    \item \textbf{Dependencias:} RF-04, RF-05.
    \item \textbf{Fuentes:} sección ``Validaciones cruzadas'' del reto.
\end{itemize}

\subsection*{Dictamen}

\textbf{RF-09 Nivel de confianza}
\begin{itemize}
    \item \textbf{Descripción:} el sistema deberá calcular y mostrar un nivel de confianza del análisis (global y por hallazgo).
    \item \textbf{Razones:} indica al equipo legal dónde conviene profundizar la revisión.
    \item \textbf{Dependencias:} RF-06, RF-08.
    \item \textbf{Fuentes:} entregables del reto.
\end{itemize}

\textbf{RF-10 Recomendación inicial}
\begin{itemize}
    \item \textbf{Descripción:} el sistema deberá emitir una recomendación inicial: Cumple, Requiere revisión o No cumple.
    \item \textbf{Razones:} permite priorizar expedientes y detectar pronto los que claramente no cumplen.
    \item \textbf{Dependencias:} RF-06, RF-07, RF-08, RF-09.
    \item \textbf{Fuentes:} planteamiento del Reto 1.
\end{itemize}

\textbf{RF-11 Explicación de la recomendación}
\begin{itemize}
    \item \textbf{Descripción:} el sistema deberá explicar con claridad las razones de la recomendación, citando los hallazgos y su evidencia.
    \item \textbf{Razones:} facilita la revisión del equipo legal sin abrir el expediente.
    \item \textbf{Dependencias:} RF-10, RF-05.
    \item \textbf{Fuentes:} planteamiento del Reto 1.
\end{itemize}

\subsection*{Presentación de resultados}

\textbf{RF-12 Tablero de resultados del expediente}
\begin{itemize}
    \item \textbf{Descripción:} el sistema deberá presentar documentos recibidos y faltantes, datos extraídos, criterios y validaciones realizadas, requisitos cumplidos y no cumplidos, inconsistencias, evidencia, nivel de confianza y recomendación.
    \item \textbf{Razones:} es el entregable principal exigido por el reto.
    \item \textbf{Dependencias:} RF-02 a RF-11.
    \item \textbf{Fuentes:} sección ``Qué entregas''.
\end{itemize}

\textbf{RF-13 Bitácora de ejecución de agentes}
\begin{itemize}
    \item \textbf{Descripción:} el sistema deberá registrar qué agente produjo cada resultado y en qué orden, para demostrar la colaboración entre agentes.
    \item \textbf{Razones:} sustenta el criterio de arquitectura multiagente y facilita la demostración en vivo.
    \item \textbf{Dependencias:} RF-02 a RF-11.
    \item \textbf{Fuentes:} criterio de Arquitectura multiagente y prototipo (25\%).
\end{itemize}

% ---------------------------------------------------------
\section{Especificaciones de los requerimientos funcionales}
% ---------------------------------------------------------

\textbf{RF-01 Carga de expediente}
\begin{itemize}
    \item \textbf{RF-01.1} El sistema deberá aceptar archivos PDF individuales o varios a la vez (incluido un .zip con PDF).
    \item \textbf{RF-01.2} El sistema deberá validar que los archivos sean PDF legibles y rechazar los corruptos con un mensaje claro.
    \item \textbf{RF-01.3} El sistema deberá agrupar los documentos en un expediente con identificador único.
    \item \textbf{RF-01.4} El sistema deberá iniciar el análisis automáticamente tras la carga, sin pasos manuales intermedios.
\end{itemize}

\textbf{RF-02 Clasificación}
\begin{itemize}
    \item \textbf{RF-02.1} El sistema deberá extraer el texto de cada PDF; si es escaneado, deberá aplicar OCR.
    \item \textbf{RF-02.2} El sistema deberá asignar un tipo documental y un grado de certeza de la clasificación.
    \item \textbf{RF-02.3} El sistema deberá marcar como ``no reconocido'' y enviar a revisión los documentos con certeza baja.
    \item \textbf{RF-02.4} El sistema deberá permitir que el usuario corrija manualmente la clasificación.
\end{itemize}

\textbf{RF-03 Faltantes}
\begin{itemize}
    \item \textbf{RF-03.1} El sistema deberá comparar los tipos presentes contra los cuatro tipos mínimos.
    \item \textbf{RF-03.2} El sistema deberá listar los tipos faltantes y los documentos duplicados o sobrantes.
\end{itemize}

\textbf{RF-04 Extracción}
\begin{itemize}
    \item \textbf{RF-04.1} El sistema deberá usar un esquema de campos por tipo documental con salida estructurada.
    \item \textbf{RF-04.2} El sistema deberá devolver ``no encontrado'' cuando un campo no exista, sin inventar valores.
    \item \textbf{RF-04.3} El sistema deberá normalizar datos para su comparación (direcciones, nombres de propietarios, unidades de superficie, fechas).
\end{itemize}

\textbf{RF-05 Evidencia}
\begin{itemize}
    \item \textbf{RF-05.1} El sistema deberá guardar para cada dato: documento, número de página y fragmento literal.
    \item \textbf{RF-05.2} Un agente verificador deberá comprobar que el fragmento citado exista en el documento original; si no existe, el dato se descarta o se marca como no verificado.
    \item \textbf{RF-05.3} El sistema deberá permitir abrir el documento en la página citada desde el tablero.
\end{itemize}

\textbf{RF-06 Validación por reglas}
\begin{itemize}
    \item \textbf{RF-06.1} El sistema deberá cargar las reglas por tipo documental desde configuración externa (no escritas en el código).
    \item \textbf{RF-06.2} El sistema deberá marcar cada regla como cumple, no cumple o indeterminada, con su evidencia.
    \item \textbf{RF-06.3} En Uso de Suelo, el sistema deberá determinar si autoriza el uso comercial requerido e identificar restricciones y condicionantes.
    \item \textbf{RF-06.4} El sistema deberá verificar fecha de emisión y vigencia cuando el tipo documental lo requiera.
\end{itemize}

\textbf{RF-07 Incompletitud y riesgos}
\begin{itemize}
    \item \textbf{RF-07.1} El sistema deberá reportar campos obligatorios no encontrados por documento.
    \item \textbf{RF-07.2} El sistema deberá reportar documentos vencidos, ilegibles o con baja calidad de OCR.
    \item \textbf{RF-07.3} El sistema deberá clasificar cada hallazgo por severidad (crítico, relevante, informativo).
\end{itemize}

\textbf{RF-08 Validación cruzada}
\begin{itemize}
    \item \textbf{RF-08.1} El sistema deberá comparar dirección, propietario o titular, superficie, número oficial e identificación del predio entre todos los documentos aplicables.
    \item \textbf{RF-08.2} El sistema deberá aplicar comparación flexible para diferencias de formato (abreviaturas, acentos, mayúsculas) y distinguirlas de discrepancias reales.
    \item \textbf{RF-08.3} El sistema deberá aplicar tolerancias de superficie [por confirmar con las reglas oficiales].
    \item \textbf{RF-08.4} El sistema deberá señalar posible pertenencia a inmuebles distintos cuando varios identificadores no coincidan.
    \item \textbf{RF-08.5} Para cada discrepancia, el sistema deberá mostrar documentos involucrados, valor en cada uno y evidencia.
\end{itemize}

\textbf{RF-09 Confianza}
\begin{itemize}
    \item \textbf{RF-09.1} El sistema deberá calcular la confianza a partir de la calidad de OCR, la certeza de clasificación, la verificación de evidencia y la concordancia entre agentes.
    \item \textbf{RF-09.2} El sistema deberá expresar la confianza en una escala comprensible (alta, media, baja) con su justificación.
\end{itemize}

\textbf{RF-10 Recomendación}
\begin{itemize}
    \item \textbf{RF-10.1} El sistema deberá emitir No cumple solo cuando exista incumplimiento claro sustentado en evidencia (por ejemplo, uso de suelo no compatible con el uso comercial).
    \item \textbf{RF-10.2} El sistema deberá emitir Requiere revisión cuando existan discrepancias, faltantes subsanables, evidencia insuficiente o baja confianza.
    \item \textbf{RF-10.3} El sistema deberá emitir Cumple solo cuando todos los requisitos se cumplan, no existan inconsistencias críticas y la confianza sea suficiente.
\end{itemize}

\textbf{RF-11 Explicación}
\begin{itemize}
    \item \textbf{RF-11.1} El sistema deberá generar un resumen en lenguaje claro con las razones principales de la recomendación.
    \item \textbf{RF-11.2} Cada razón deberá enlazar a su hallazgo y evidencia.
\end{itemize}

\textbf{RF-12 Tablero}
\begin{itemize}
    \item \textbf{RF-12.1} El sistema deberá mostrar el progreso del análisis en tiempo real.
    \item \textbf{RF-12.2} El sistema deberá permitir filtrar hallazgos por severidad, documento o tipo (faltante, inconsistencia, incumplimiento).
    \item \textbf{RF-12.3} El sistema deberá permitir exportar el dictamen a PDF o JSON.
\end{itemize}

\textbf{RF-13 Bitácora}
\begin{itemize}
    \item \textbf{RF-13.1} El sistema deberá registrar entradas, salidas y tiempos de cada agente por expediente.
    \item \textbf{RF-13.2} El sistema deberá mostrar el flujo de agentes de la corrida en la interfaz para la demostración.
\end{itemize}

% ---------------------------------------------------------
\section{Formalización de los requerimientos no funcionales}
% ---------------------------------------------------------

Las metas numéricas son propuestas del equipo y se ajustarán con los expedientes de prueba y las reglas del comité.

\textbf{RNF-01 Precisión}
\begin{itemize}
    \item \textbf{Descripción:} la extracción, la validación contra reglas y las validaciones cruzadas deben ser correctas, sin descartar expedientes que sí podían continuar.
    \item \textbf{Métrica:} coincidencia con el dictamen de referencia en los expedientes de prueba [meta: $\ge$ 85\%]; tasa de falsos ``No cumple'' sobre expedientes viables lo más cercana a 0.
    \item \textbf{Criterio de aceptación:} la corrida sobre los expedientes de prueba se compara contra el dictamen de referencia y se documentan los desacuerdos.
\end{itemize}

\textbf{RNF-02 Explicabilidad y trazabilidad}
\begin{itemize}
    \item \textbf{Descripción:} cada hallazgo debe citar la evidencia del documento original.
    \item \textbf{Métrica:} 100\% de los hallazgos con documento, página y fragmento; 0 citas no verificadas en el resultado final.
    \item \textbf{Criterio de aceptación:} un revisor entiende la recomendación sin abrir el expediente.
\end{itemize}

\textbf{RNF-03 Arquitectura multiagente}
\begin{itemize}
    \item \textbf{Descripción:} el problema debe dividirse en agentes especializados que contrastan y validan resultados entre sí, sobre LangChain.
    \item \textbf{Métrica:} al menos un agente de verificación independiente de los agentes que extraen y validan; flujo completo de punta a punta.
    \item \textbf{Criterio de aceptación:} demostración en vivo con un expediente cargado durante la sesión.
\end{itemize}

\textbf{RNF-04 Rendimiento}
\begin{itemize}
    \item \textbf{Métrica:} análisis de un expediente completo en $\le$ 5 minutos [meta propuesta]; progreso visible durante la ejecución.
    \item \textbf{Criterio de aceptación:} la demostración en vivo corre sin esperas que afecten el pitch.
\end{itemize}

\textbf{RNF-05 Confidencialidad y seguridad}
\begin{itemize}
    \item \textbf{Descripción:} los documentos del reto se usan solo dentro del reto y no se suben a otros servicios ni se comparten fuera del equipo.
    \item \textbf{Métrica:} cero envíos de documentos a servicios no autorizados; registros locales sin contenido sensible innecesario.
    \item \textbf{Criterio de aceptación:} antes de usar un proveedor de modelos de lenguaje o de OCR, se confirma con el comité que está permitido; si no, se usan alternativas locales. [por confirmar]
\end{itemize}

\textbf{RNF-06 Robustez de ingesta}
\begin{itemize}
    \item \textbf{Descripción:} el sistema debe manejar PDF digitales y escaneados, con calidad variable.
    \item \textbf{Métrica:} tasa de documentos procesados sin error; los fallos se reportan, no detienen el expediente.
\end{itemize}

\textbf{RNF-07 Usabilidad}
\begin{itemize}
    \item \textbf{Descripción:} el equipo legal debe poder cargar un expediente y entender el resultado sin capacitación.
    \item \textbf{Métrica:} flujo de carga a dictamen en $\le$ 3 acciones del usuario; lenguaje no técnico en el tablero.
\end{itemize}

\textbf{RNF-08 Mantenibilidad y extensibilidad}
\begin{itemize}
    \item \textbf{Descripción:} agregar un tipo documental o una regla nueva no debe requerir reescribir el flujo.
    \item \textbf{Métrica:} reglas y esquemas en configuración; documentación de la arquitectura y del diagrama de agentes; al menos 2 refactorizaciones justificadas antes de la entrega.
\end{itemize}

\textbf{RNF-09 Reproducibilidad}
\begin{itemize}
    \item \textbf{Descripción:} la misma entrada debe producir resultados consistentes.
    \item \textbf{Métrica:} temperatura baja en los modelos, salidas estructuradas y versiones fijas de prompts y reglas.
\end{itemize}

% ---------------------------------------------------------
\section{Reglas de negocio}
% ---------------------------------------------------------

Son restricciones del dominio que el sistema debe aplicar automáticamente. Las siguientes provienen del planteamiento del reto; las reglas por tipo documental se completarán con el material oficial del comité.

\begin{enumerate}
    \item Un expediente debe contener como mínimo cuatro tipos documentales: Escrituras, Predial, Alineamiento y Número Oficial, y Uso de Suelo.
    \item La recomendación inicial solo puede ser \textit{Cumple}, \textit{Requiere revisión} o \textit{No cumple}.
    \item Toda conclusión debe estar sustentada en evidencia de los documentos originales; sin evidencia no hay hallazgo.
    \item Ante información insuficiente o ambigua, la recomendación debe ser \textit{Requiere revisión}; el sistema no descarta expedientes que podían continuar.
    \item El uso de suelo autorizado debe ser compatible con el uso comercial requerido; las restricciones o condicionantes deben reportarse.
    \item Los documentos con fecha de emisión y vigencia (Alineamiento y Número Oficial, Uso de Suelo) deben estar vigentes cuando corresponda.
    \item Deben coincidir entre documentos del mismo inmueble: dirección, propietario o titular, superficie (dentro de tolerancia), número oficial e identificación del predio; toda diferencia relevante se reporta.
    \item Los documentos del reto se usan solo dentro del reto.
    \item Ante cualquier diferencia entre este documento y los documentos oficiales del comité, mandan los oficiales.
    \item \textbf{[Por confirmar]} Tolerancia de superficie, criterios exactos de vigencia, lista de autoridades emisoras válidas y uso de suelo considerado compatible con tienda o centro de distribución.
\end{enumerate}

% ---------------------------------------------------------
\section{Historias de usuario}
% ---------------------------------------------------------

\textbf{Actores:} analista legal (usuario principal) y coordinador del equipo legal.

\textbf{HU-01 Cargar expediente.} Como analista legal, quiero cargar los documentos de un inmueble, para iniciar la revisión inicial automática.
\begin{itemize}
    \item Dado que tengo los PDF de un inmueble, cuando los cargo, entonces el sistema crea el expediente e inicia el análisis.
    \item Dado que un archivo no es un PDF válido, cuando lo cargo, entonces el sistema me avisa cuál falló y continúa con los demás.
\end{itemize}

\textbf{HU-02 Clasificación automática.} Como analista legal, quiero que el sistema identifique el tipo de cada documento, para no clasificar manualmente.
\begin{itemize}
    \item Dado un documento cargado, cuando termina el análisis, entonces veo su tipo y la certeza de la clasificación.
    \item Dado un documento no reconocido, entonces el sistema lo marca para mi revisión.
\end{itemize}

\textbf{HU-03 Ver faltantes.} Como analista legal, quiero ver qué documentos faltan, para solicitarlos antes de seguir.
\begin{itemize}
    \item Dado un expediente sin un tipo mínimo, cuando abro el resultado, entonces veo el tipo faltante señalado.
\end{itemize}

\textbf{HU-04 Ver datos extraídos.} Como analista legal, quiero ver los datos relevantes de cada documento, para revisarlos sin leer todo el PDF.
\begin{itemize}
    \item Dado un documento analizado, cuando lo abro, entonces veo cada campo con su valor y su evidencia.
\end{itemize}

\textbf{HU-05 Ver validaciones por reglas.} Como analista legal, quiero ver qué requisitos se cumplen y cuáles no, para entender el estado del expediente.
\begin{itemize}
    \item Dado un expediente analizado, cuando consulto las validaciones, entonces veo cada regla con su resultado y su evidencia.
\end{itemize}

\textbf{HU-06 Ver inconsistencias entre documentos.} Como analista legal, quiero ver las discrepancias entre documentos, para detectar posibles errores o inmuebles distintos.
\begin{itemize}
    \item Dado que dos documentos difieren en un dato, cuando veo el análisis, entonces se muestra el dato, el valor en cada documento y la evidencia.
\end{itemize}

\textbf{HU-07 Verificar evidencia.} Como analista legal, quiero abrir el fragmento que sustenta un hallazgo, para confirmarlo rápidamente.
\begin{itemize}
    \item Dado un hallazgo, cuando selecciono su evidencia, entonces veo el documento en la página y el fragmento citado.
\end{itemize}

\textbf{HU-08 Recibir recomendación.} Como analista legal, quiero una recomendación inicial con sus razones y su confianza, para priorizar mi trabajo.
\begin{itemize}
    \item Dado un análisis completo, cuando termina, entonces veo Cumple, Requiere revisión o No cumple, con razones y nivel de confianza.
    \item Dado que existe duda, entonces la recomendación es Requiere revisión y no No cumple.
\end{itemize}

\textbf{HU-09 Corregir el resultado.} Como analista legal, quiero corregir una clasificación o un dato extraído y volver a ejecutar la validación, para obtener un dictamen preciso.
\begin{itemize}
    \item Dado un dato incorrecto, cuando lo corrijo, entonces el sistema recalcula validaciones y recomendación.
\end{itemize}

\textbf{HU-10 Exportar dictamen.} Como coordinador, quiero exportar el dictamen del expediente, para compartirlo con el equipo legal.
\begin{itemize}
    \item Dado un análisis terminado, cuando lo exporto, entonces obtengo un documento con hallazgos, evidencia y recomendación.
\end{itemize}

\textbf{HU-11 Priorizar expedientes.} Como coordinador, quiero ver varios expedientes con su recomendación, para decidir cuáles revisar primero.
\begin{itemize}
    \item Dado que hay varios expedientes analizados, cuando abro la lista, entonces los veo con su recomendación y confianza.
\end{itemize}

% ---------------------------------------------------------
\section{Actores y casos de uso}
% ---------------------------------------------------------

\subsection*{Actores}
\begin{itemize}
    \item \textbf{Analista legal:} carga expedientes, consulta resultados, revisa evidencia, corrige datos y exporta el dictamen.
    \item \textbf{Coordinador del equipo legal:} consulta la lista de expedientes y prioriza la revisión.
    \item \textbf{Administrador técnico:} mantiene la configuración de reglas y esquemas por tipo documental.
    \item \textbf{Servicios externos o locales (modelo de lenguaje, OCR):} actores secundarios que dan soporte a los agentes; su uso está sujeto a RNF-05.
\end{itemize}

\subsection*{Casos de uso principales}
\begin{itemize}
    \item Cargar expediente (incluye: validar archivos, crear expediente).
    \item Clasificar documentos (incluye: extraer texto/OCR, asignar tipo).
    \item Extraer datos (incluye: aplicar esquema por tipo, registrar evidencia).
    \item Validar contra reglas (incluye: cargar reglas, evaluar cada regla).
    \item Validar entre documentos (incluye: normalizar datos, comparar campos, detectar posibles inmuebles distintos).
    \item Emitir recomendación (incluye: calcular confianza, generar explicación).
    \item Consultar resultados y evidencia (extiende: corregir dato y recalcular).
    \item Exportar dictamen.
    \item Configurar reglas y esquemas.
\end{itemize}

\noindent\textit{Nota: Los diagramas de casos de uso, de actividad (análisis de expediente, validación cruzada, dictamen) y de secuencia (flujo entre agentes) se anexarán a este documento.}

% ---------------------------------------------------------
\section{Arquitectura multiagente}
% ---------------------------------------------------------

El problema se divide en agentes especializados construidos con LangChain, coordinados por un orquestador que mantiene un estado compartido del expediente. La coordinación puede implementarse como grafo de agentes con LangGraph, del ecosistema de LangChain \texttt{[por confirmar que cumple el stack obligatorio]}.

\subsection*{Agentes}
\begin{itemize}
    \item \textbf{Agente de ingesta:} lee los PDF, extrae texto, aplica OCR si es necesario, segmenta con referencias de documento y página.
    \item \textbf{Agente clasificador:} identifica el tipo documental y su certeza.
    \item \textbf{Agentes extractores (uno por tipo documental):} extraen los campos definidos con salida estructurada y evidencia.
    \item \textbf{Agentes validadores (uno por tipo documental):} evalúan las reglas y criterios del tipo.
    \item \textbf{Agente de validación cruzada:} compara dirección, propietario, superficie, número oficial e identificación del predio entre documentos.
    \item \textbf{Agente verificador (crítico):} revisa de forma independiente que cada dato y cada cita existan en el documento original y que los hallazgos sean consistentes; rechaza o degrada lo no verificado.
    \item \textbf{Agente dictaminador:} integra hallazgos, calcula confianza, emite la recomendación y redacta la explicación.
    \item \textbf{Orquestador:} decide el orden, reintenta o escala a revisión humana cuando hay conflicto entre agentes.
\end{itemize}

\subsection*{Flujo}
\begin{enumerate}
    \item Ingesta de los PDF del expediente.
    \item Clasificación de cada documento y detección de faltantes.
    \item Extracción por tipo, con evidencia.
    \item Verificación de evidencia por el agente verificador.
    \item Validación por reglas por tipo.
    \item Validación cruzada entre documentos.
    \item Contraste entre agentes: si extractor y verificador discrepan, el dato queda como no verificado y baja la confianza.
    \item Dictamen: recomendación, confianza y explicación.
    \item Presentación en el tablero y bitácora de agentes.
\end{enumerate}

\subsection*{Criterios del dictamen}
\begin{itemize}
    \item \textbf{No cumple:} incumplimiento claro y sustentado en evidencia verificada.
    \item \textbf{Cumple:} todos los requisitos cumplidos, sin inconsistencias críticas y con confianza suficiente.
    \item \textbf{Requiere revisión:} cualquier otro caso, incluyendo evidencia insuficiente, discrepancias subsanables y documentos faltantes.
\end{itemize}

% ---------------------------------------------------------
\section{Patrones arquitectónicos y de diseño}
% ---------------------------------------------------------

\subsection*{Patrones arquitectónicos}
\begin{itemize}
    \item \textbf{Orquestador de agentes (supervisor):} un agente coordina a los especialistas y decide el siguiente paso según el estado.
    \item \textbf{Pipeline por etapas:} ingesta, clasificación, extracción, validación, validación cruzada, dictamen, con interfaces claras entre etapas.
    \item \textbf{Estado compartido (pizarra):} el estado del expediente concentra documentos, datos, evidencia y hallazgos para que los agentes lo consulten y lo enriquezcan.
    \item \textbf{Arquitectura por capas:} presentación (interfaz), servicios (casos de uso), dominio (reglas) y datos (persistencia).
    \item \textbf{Salidas estructuradas:} los agentes devuelven objetos con esquema definido, no texto libre, para validar y trazar.
\end{itemize}

\subsection*{Patrones de diseño}
\begin{itemize}
    \item \textbf{Strategy (reglas por tipo documental):} cada regla o conjunto de reglas se encapsula en una clase intercambiable, de modo que agregar una regla no modifica las existentes.
    \item \textbf{Factory Method (agentes por tipo):} crea el extractor y el validador adecuados según el tipo clasificado.
    \item \textbf{Template Method (flujo de extracción):} define el esqueleto común (leer, extraer, citar, validar) y deja variar los pasos por tipo.
    \item \textbf{Facade (análisis de expediente):} una clase de servicio ofrece un único punto de entrada a todo el flujo para la interfaz.
    \item \textbf{Adapter (modelo de lenguaje y OCR):} aísla al proveedor concreto para poder cambiarlo si el comité restringe servicios.
    \item \textbf{Observer (progreso y bitácora):} los agentes emiten eventos que alimentan la barra de progreso y el registro de ejecución.
    \item \textbf{Specification (filtros y condiciones):} composición de condiciones para filtrar hallazgos y expedientes.
\end{itemize}

% ---------------------------------------------------------
\section{Mapa de navegación}
% ---------------------------------------------------------

\begin{itemize}
    \item \textbf{Inicio}
    \begin{itemize}
        \item \textbf{Cargar expediente}
        \begin{itemize}
            \item Selección de archivos
            \item Progreso del análisis (agentes en ejecución)
        \end{itemize}
        \item \textbf{Lista de expedientes}
        \begin{itemize}
            \item Recomendación y confianza por expediente
            \item Filtros por recomendación y fecha
        \end{itemize}
        \item \textbf{Detalle de expediente}
        \begin{itemize}
            \item Resumen y recomendación (razones y confianza)
            \item Documentos recibidos y faltantes
            \item Datos extraídos por documento
            \item Validaciones por reglas (cumplidas y no cumplidas)
            \item Inconsistencias entre documentos
            \item Visor de evidencia (documento, página, fragmento)
            \item Bitácora de agentes
            \item Exportar dictamen
        \end{itemize}
        \item \textbf{Configuración (administrador)}
        \begin{itemize}
            \item Reglas por tipo documental
            \item Esquemas de campos
        \end{itemize}
    \end{itemize}
\end{itemize}

% ---------------------------------------------------------
\section{Aplicación del ciclo de vida del software en el proyecto}
% ---------------------------------------------------------

El equipo sigue las etapas de forma iterativa y adaptada a los tiempos del hackatón: sprints cortos, revisiones frecuentes del equipo y tablero de tareas.

\begin{itemize}
    \item \textbf{Planificación.} Se definió el alcance a partir del planteamiento del Reto 1, se asignaron roles y se estableció la forma de trabajo. \textit{Pendiente:} calendario exacto según las fechas del hackatón.
    \item \textbf{Recolección y análisis de requerimientos.} Se documentaron los requerimientos funcionales y no funcionales, las reglas de negocio y las historias de usuario con criterios de aceptación. \textit{Pendiente:} incorporar las reglas y criterios definitivos que entrega el comité al arranque.
    \item \textbf{Diseño del sistema.} Se propuso la arquitectura multiagente, el flujo de agentes, los patrones y el mapa de navegación. \textit{Pendiente:} decisión final sobre la herramienta de orquestación y la pila tecnológica (backend, interfaz, OCR, proveedor de modelos) conforme a las restricciones de confidencialidad.
    \item \textbf{Implementación.} Se desarrollará por incrementos verticales: primero un expediente de punta a punta con un tipo documental, después los demás tipos, la validación cruzada y el dictamen. Se usará control de versiones y al menos dos refactorizaciones documentadas.
    \item \textbf{Pruebas.} Se diseñarán pruebas de unidad (extracción y reglas), integración (flujo entre agentes) y sistema (expedientes de prueba contra el dictamen de referencia), con atención a los casos donde un expediente viable podría descartarse por error.
    \item \textbf{Despliegue.} Se prevé una ejecución local o en un entorno controlado que permita la demostración en vivo con un expediente cargado durante la sesión, sin enviar documentos a servicios no autorizados.
    \item \textbf{Mantenimiento y evolución.} El diseño por configuración (reglas y esquemas externos) busca que se puedan agregar tipos documentales y reglas sin reescribir el flujo.
    \item \textbf{Retiro o reemplazo.} No se contempla por el momento.
\end{itemize}

\subsection*{Entregables y pitch}
\begin{itemize}
    \item Prototipo funcional que ingresa un expediente, ejecuta el análisis automáticamente y presenta todos los resultados descritos en RF-12.
    \item Diagrama de agentes.
    \item Pitch de 3 minutos: problema, cómo funciona la solución, cómo llega a la recomendación y el valor para el equipo legal.
\end{itemize}

\end{document}